\hdr{§0}{先看这一页：这份手册是什么、今晚怎么用}

\callout{\textbf{一句话定位。}明天这场考的不是“多学了多少章”，而是\emph{把学过的东西在 150 分钟里准确写出来}的能力。非数学类初赛的全部内容是高等数学：极限与连续、一元微分、一元积分、向量代数与空间解析几何、多元微分、重积分、曲线与曲面积分、无穷级数、常微分方程。你现在有微积分基础，缺的主要是\emph{下册}（多元、级数、曲线曲面、方程）与\emph{竞赛式套路}。这份手册把这两块按“能读懂”的标准从零补上，并把上册内容重新整理成可执行的解题动作。}

\subhdr{一、范围与卷面}

\begin{center}
\tabtop
\begin{tabular}{@{}p{3.4cm}p{9.6cm}@{}}
\thead{项目} & \thead{说明}\\
\hline
考试 & 全国大学生数学竞赛\textbf{预赛}（初赛），2026-11-14（周六）09:00--11:30，150 分钟\\
组别 & 你是物理专业，对应\textbf{非数学类}；命题范围就是高等数学上下册，不考数学分析、高等代数\\
卷面 & 近年非数学类初赛满分 100 分，常见结构为填空题 5 道（每题 5 分，共 25 分）＋解答题 5 道（每题 15 分，共 75 分）。\textbf{实际结构以你拿到的试卷为准}，填空做完先数一遍题号\\
难度 & 比期末难在：(1)计算量更大；(2)证明题要自己构造；(3)喜欢把两章内容接在一起考\\
\end{tabular}
\tabbot
\end{center}

\warnbox{\textbf{两条边界，先说清楚。}
（1）上述范围来自竞赛公开大纲与历年卷面，你的准考证／校内通知若另写范围，以那份为准。
（2）这份手册按\emph{非数学类}组织。若你实际报的是数学类，那需要另补实数完备性、一致收敛、多元微分的严格理论以及高等代数，本手册的极限、级数、多元部分仍是共同底座，但不能覆盖数学类的缺口——明天前补不完，就先抓本手册里标 \ding{72} 的计算题。}

\subhdr{二、优先级地图：时间该花在哪}

下表按“考频 × 提分速度”排序。\ding{72}越多越该先看；“读”指读懂讲解，“练”指动笔复现。

\begin{center}
\small
\tabtop
\begin{tabular}{@{}clcp{6.6cm}@{}}
\thead{优先级} & \thead{模块} & \thead{典型分值} & \thead{明天最可能怎么考}\\
\hline
\ding{72}\ding{72}\ding{72}\ding{72}\ding{72} & §3 一元积分 & 15--30 & 定积分技巧题、对称性与区间再现、变上限积分、积分不等式\\
\ding{72}\ding{72}\ding{72}\ding{72}\ding{72} & §2 一元微分 & 15--30 & 中值定理证明（构造辅助函数）、不等式证明、高阶导\\
\ding{72}\ding{72}\ding{72}\ding{72}\ding{72} & §1 极限 & 5--15 & 填空必有一道：等价无穷小／泰勒／$1^\infty$／$n$ 项和化定积分\\
\ding{72}\ding{72}\ding{72}\ding{72} & §6 曲线曲面积分 & 15 & 格林公式（补线／挖洞）、高斯公式（补面）\\
\ding{72}\ding{72}\ding{72}\ding{72} & §7 级数 & 15 & 判敛、幂级数收敛域、和函数\\
\ding{72}\ding{72}\ding{72}\ding{72} & §8 常微分方程 & 15 & 一阶线性／伯努利、二阶常系数、欧拉方程、由积分方程反推\\
\ding{72}\ding{72}\ding{72} & §5 重积分 & 15 & 极坐标、交换次序、对称性\\
\ding{72}\ding{72}\ding{72} & §4 多元微分 & 5--15 & 偏导与全微分概念、条件极值（拉格朗日）\\
\ding{72}\ding{72} & §9 向量与空间几何 & 5 & 填空：平面／直线方程、距离、夹角、旋转曲面\\
\ding{72}\ding{72} & §10 套路专题 & 贯穿 & 上面各章的“通用招式”汇总，最后读\\
\ding{72} & §11 考场执行 & — & 明早出门前 10 分钟读\\
\end{tabular}
\tabbot
\end{center}

\subhdr{三、三种时间预算，照自己的情况选一条}

\callout{\textbf{A. 今晚有 3 小时（推荐）。}
\begin{enumerate}[leftmargin=1.6em,itemsep=1pt]
\item 20:00--20:50 读 §1、§2：把“极限四招”和“中值定理构造七模板”抄到一张纸上；
\item 20:50--21:40 读 §3、§4、§5：重点是积分技巧表、极坐标与换序、拉格朗日乘数；
\item 21:40--22:30 读 §6、§7、§8：格林／高斯各背一个标准流程，级数判敛表，ODE 五种方程；
\item 22:30--22:50 读 §10 与附录 A，只记模板名，不做题；
\item 睡前不再碰新内容。明早 7:30--8:20 只翻附录 A 速查表与 §11 清单。
\end{enumerate}}

\callout{\textbf{B. 今晚只有 90 分钟。}只读 §3（积分）、§6（格林／高斯）、§7（级数判敛与幂级数）、§8（ODE 五种方程），这四块合计约 60--70 分，且套路化最强。§1、§2 只看附录 A 的公式表。\textbf{放弃}§9 与 §10 的细读。}

\callout{\textbf{C. 明早只有 40 分钟。}只看三样：附录 A 速查表、§11 的失分点清单、§8 的“二阶常系数特解设法”表。别的都不要翻，翻了只会慌。}

\subhdr{四、这份手册的记号约定}

\begin{itemize}[leftmargin=1.6em,itemsep=1pt]
\item 标 \ding{72} 的是“明天很可能出现、且能在 10 分钟内学会”的条目；标「低」的是低频内容，时间不够可跳过。
\item “识别”指题目长什么样；“动作”指第一步该写什么。竞赛题的差距通常不在会不会，而在\emph{第一步有没有写对}。
\item 所有例题都是\textbf{教学小例}：用来演示方法，不是给你做的练习。每道都给了完整过程，照着走一遍比看懂更重要。
\item 附录 B 有 10 道自测题（答案在最后），明早用来检查手感，不占用今晚时间。
\end{itemize}

\warnbox{\textbf{关于“掌握”的实话。}读完这份手册，你能做到的是：认得每一类题的入口、写得出标准步骤、不犯结构性错误。这足以把 100 分的卷子拿到 60--75 分。真正的稳定发挥来自动笔，而今晚的时间只够你\emph{每类题动笔一道}——所以 §0 的三条路线里，每一段都留了 5--10 分钟给“合上书自己写一遍关键式”。跳过这一步，明天看到题会觉得眼熟但写不出来。}
